\section{PX4 lockstep integration}
\label{sec:px4-bridge}

PX4 is treated as a deterministic discrete-time component stepped by integer microseconds. The simulation engine never advances PX4 by wall-clock time; instead it calls the lockstep step function with an exact \code{UInt64} timestamp at each autopilot boundary.

\subsection{Autopilot step interface}

The canonical interface is \code{Autopilots.autopilot_step} (\code{src/sim/Autopilots.jl}). On each autopilot tick, the engine provides:
\begin{itemize}[leftmargin=*,itemsep=2pt]
  \item estimated state: $(\hat{\vect{p}}_{ned}, \hat{\vect{v}}_{ned}, \hat{\qbn}, \hat{\vect{\omega}}_{body})$,
  \item high-level command \code{AutopilotCommand} (arm/mission/RTL),
  \item a best-effort \code{landed} flag, and
  \item per-battery \code{BatteryStatus} records computed by the plant outputs protocol.
\end{itemize}

\paragraph{State-to-global conversion.}
PX4 expects global coordinates for some topics. The implementation converts local NED to an approximate latitude/longitude/altitude using a fixed Earth radius and the configured home origin (\code{_ned\_to\_lla}). Altitude is computed as $h_{msl}=h_{home}-p_{ned,z}$. This approximation is intended for local missions (order \(\sim\)\SI{10}{km} from home) and uses a fixed \(\cos(\varphi_0)\) longitude scale; see \cref{sec:design-choices}.

\paragraph{Mode selection.}
The Julia bridge forwards the high-level command packet to the C runtime (\code{set\_cmd!}); when Commander is disabled, the lockstep C runtime publishes the minimal commander‑lite topics (\code{vehicle\_status}, \code{vehicle\_control\_mode}, \code{actuator\_armed}) and drives \code{nav\_state}/\code{arming\_state} based on command inputs and \code{mission\_result}. A configurable \code{edge\_trigger} mode determines whether mission/RTL requests are treated as edges or levels.

\subsection{Multi-rate uORB injection}

PX4Lockstep uses explicit injection sources to publish uORB messages at configurable integer-microsecond periods (\code{src/sim/Autopilots/UORBInjection.jl}). The boundary contract is:
\begin{enumerate}[leftmargin=*,itemsep=2pt]
  \item queue all uORB publishes that are \emph{due} at the current \code{time\_us},
  \item call \code{step\_uorb!(handle, time\_us)} exactly once, then
  \item read back configured uORB subscriptions.
\end{enumerate}
The default injector publishes a set of ``state injection'' topics every PX4 step (period $0$), which mirrors the historical lockstep SITL behavior.

\subsection{Actuator command extraction}

After stepping PX4, the autopilot reads actuator outputs from uORB subscriptions (\code{UORBOutputs.actuator\_motors}; \code{UORBOutputs.actuator\_servos}). It then converts these into fixed-size \code{ActuatorCommand} arrays and publishes them into the runtime bus (\code{src/sim/Sources/Autopilot.jl}).

\paragraph{Actuator range semantics.}
The bridge uses normalized uORB outputs: \code{actuator\_motors} is treated as a duty-like command in \([0,1]\), and \code{actuator\_servos} as a signed deflection in \([-1,1]\). These are lockstep ABI conventions (not PWM). If a different uORB topic or scaling is required, the mapping should be changed in the bridge configuration and documented alongside the uORB interface.

\subsection{Derived PX4 parameterization: control allocator geometry}

When enabled in the aircraft spec (\code{PX4Spec.derive\_ca\_params}), the build layer derives \code{CA\_*} control allocator parameters from the multirotor geometry (\code{src/sim/Aircraft/Build.jl}). For each rotor $i$ with body position $\vect{r}_{b,i}$ and propulsor axis $\vect{a}_{b,i}$, PX4 expects a \emph{thrust direction} axis; the builder therefore provides $-\vect{a}_{b,i}$ (recall \cref{sec:axis-convention}). The yaw moment coefficient parameter \code{CA\_ROTOR<i>\_KM} is set to \texttt{km\_m} times the rotor reaction torque sign pattern \code{rotor\_dir[i]}.

\paragraph{Caveat.}
This parameter derivation keeps PX4-side allocation consistent with the Julia-side plant geometry, but it does not configure PX4 mixers or select an airframe; those remain PX4 configuration responsibilities.


\subsection{TOML configuration: PX4 bridge, lockstep rates, and uORB interface}
\label{sec:px4-toml}

The PX4 integration is configured under \code{[px4]} (\code{PX4Spec} parsed in \code{src/sim/Aircraft/TOMLIO.jl}). The same table also contains the lockstep module scheduling rates (\code{[px4.lockstep]}) and the uORB publish/subscribe interface definition (\code{[px4.uorb]}).

\paragraph{Example (Iris defaults).}
\begin{lstlisting}[language=toml]
[px4]
libpath = "../../../../build/px4_sitl_lockstep/src/lib/px4_lockstep/libpx4_lockstep.so"
mission_path = "../../src/Workflows/assets/missions/simple_mission.waypoints"
derive_ca_params = true
edge_trigger = false

# Optional lockstep-module scheduler configuration (PX4Lockstep.LockstepConfig)
[px4.lockstep]
enable_commander = 0
commander_rate_hz = 100
enable_navigator = 0
navigator_rate_hz = 20
enable_mc_pos_control = 1
mc_pos_control_rate_hz = 100
enable_mc_att_control = 1
mc_att_control_rate_hz = 250
enable_mc_rate_control = 1
mc_rate_control_rate_hz = 250
enable_control_allocator = 1
control_allocator_rate_hz = 250
dataman_use_ram = 1

# uORB bridge configuration can be extended from an asset (e.g. iris_uorb.toml)
[px4.uorb]
pubs = ["vehicle_local_position", "vehicle_attitude", "battery_status"]
subs = ["actuator_outputs"]

# Optional PX4 parameter overrides (applied at init)
[[px4.params]]
name  = "SYS_AUTOSTART"
value = 4001

[[px4.params]]
name  = "COM_ARM_CHK"
value = 0
\end{lstlisting}

\paragraph{Parameter mapping.}
\begin{itemize}[leftmargin=*,itemsep=2pt]
  \item \code{libpath}: shared library path for \code{libpx4\_lockstep} (loaded at runtime).
  \item \code{mission\_path}: optional PX4 mission/waypoint file path. If absent, PX4 can still run in manual/offboard modes.
  \item \code{derive\_ca\_params}: when \code{true}, derive control-allocation parameters from the actuation geometry during build (see \code{Build.derive\_control\_allocator\_params}).
  \item \code{edge\_trigger}: selects whether mission/RTL requests are treated as edges (one-shot) or levels (held) in the autopilot command logic.
  \item \code{[px4.lockstep]}: configures which PX4 modules are stepped in the lockstep loop and at what rates (fields correspond one-to-one with \code{PX4Lockstep.LockstepConfig}).
  \item \code{[px4.uorb]}: declares which uORB topics are published into PX4 (\code{pubs}) and which are subscribed out of PX4 (\code{subs}); see Appendix~\ref{app:uorb-injection}.
  \item \code{[[px4.params]]}: list of \{name,value\} pairs applied to PX4 parameters at initialization.
\end{itemize}
